\mxnewpage\section*{EV3-T2 -- Blanc 2}\label{EV3-T2}
\addcontentsline{toc}{section}{EV3-T2 -- Blanc 2}
\textbf{Sujet d’entraînement}. Programme étudié des chapitres 1 à 10.\par Durée : 2 h. Sans calculatrice. Note sur 20. 1–3 avril 2027.
\subsection*{Partie 1 -- Automatismes (6 points)}
Douze questions, 0{,}5 point chacune. Une seule réponse correcte ; aucune justification demandée. Une réponse incorrecte ou absente rapporte 0, sans retrait de points.
\begin{enumerate}
\item Calculer $2+2\times3$.\par A. $11$\quad B. $7$\quad C. $12$\quad D. $8$
\item Calculer $\frac12+\frac13$.\par A. $\frac{1}{6}$\quad B. $\frac{2}{5}$\quad C. $\frac{5}{3}$\quad D. $\frac{5}{6}$
\item Calculer $\dfrac{2^2}{2^{0}}$.\par A. $7$\quad B. $4$\quad C. $2$\quad D. $1$
\item Une valeur de 50 augmente de 20 \%. Valeur finale ?\par A. $40$\quad B. $70$\quad C. $10$\quad D. $60$
\item Une valeur de 80 est obtenue après une baisse de 20 \%. Valeur initiale ?\par A. $96$\quad B. $60$\quad C. $100$\quad D. $64$
\item Hausse de 25 \% puis baisse de 20 \% : coefficient global ?\par A. $\frac{4}{5}$\quad B. $\frac{21}{20}$\quad C. $1$\quad D. $\frac{9}{10}$
\item La solution de $2x+2=0$ vaut\par A. $-1$\quad B. $4$\quad C. $-4$\quad D. $1$
\item La pente entre $(0;2)$ et $(2;8)$ vaut\par A. $3$\quad B. $2$\quad C. $6$\quad D. $-3$
\item La moyenne de $2, 4, 6$ vaut\par A. $2$\quad B. $4$\quad C. $5$\quad D. $12$
\item Un secteur de $90^\circ$ représente quelle fraction du total ?\par A. $\frac{1}{4}$\quad B. $\frac{1}{2}$\quad C. $\frac{1}{3}$\quad D. $\frac{3}{4}$
\item Une probabilité vaut $\dfrac{2}{5}$. Son complément vaut\par A. $\frac{5}{2}$\quad B. $\frac{3}{5}$\quad C. $\frac{1}{5}$\quad D. $\frac{2}{5}$
\item Pour $f(x)=x^2+2$, $f(-2)$ vaut\par A. $6$\quad B. $-2$\quad C. $4$\quad D. $8$
\end{enumerate}
\mxnewpage\subsection*{Partie 2 -- Exercices indépendants (14 points)}
Toute démarche doit être justifiée. Les résultats exacts sont attendus sauf indication contraire.
\mxneed{12}
\subsection*{Exercice 1 -- Simulation et moyenne (5 points)}
La variable aléatoire $X$ prend les valeurs 0, 2 et 6 avec les probabilités respectives $\dfrac12$, $\dfrac14$ et $\dfrac14$. On simule ensuite des échantillons de 25 réalisations indépendantes de cette variable et l’on note $m$ la moyenne observée dans un échantillon.
\begin{enumerate}
\item Calculer l’espérance $E(X)$. \hfill(1 pt)
\item Calculer $V(X)$ puis l’écart type exact $\sigma(X)$. \hfill(2 pts)
\item Écrire l’intervalle des valeurs de $m$ vérifiant $|m-E(X)|\leq\dfrac{2\sigma(X)}5$. Expliquer pourquoi une simulation sur un nombre fini d’échantillons ne garantit pas un pourcentage de couverture exact. \hfill(2 pts)
\end{enumerate}
\mxneed{12}
\subsection*{Exercice 2 -- Comparer deux modèles (5 points)}
On définit sur $\R$ les fonctions $f(x)=e^{2x}$ et $g(x)=1+2x$, puis leur différence $h=f-g$. La dernière question considère les valeurs de $f$ aux instants entiers, à partir du rang 0.
\begin{enumerate}
\item Calculer la dérivée de $h=f-g$. \hfill(1 pt)
\item Étudier h et comparer les courbes. \hfill(2 pts)
\item Pour $u_n=e^{2n}$, donner raison et somme des quatre premiers termes. \hfill(2 pts)
\end{enumerate}
\mxneed{12}
\subsection*{Exercice 3 -- Un cercle déplacé (4 points)}
Dans un repère orthonormé, on considère le cercle $\mathcal C$ d’équation $(x-1)^2+(y+1)^2=4$ et la droite $d$ d’équation $y=-1$. On nommera $A$ et $B$ leurs points d’intersection, $A$ étant celui de plus petite abscisse.
\begin{enumerate}
\item Donner centre et rayon. \hfill(1 pt)
\item Déterminer les intersections avec la droite. \hfill(1 pt)
\item Pour M sur le cercle, calculer $\vect{MA}\cdot\vect{MB}$ où A et B sont ces intersections ; justifier. \hfill(2 pts)
\end{enumerate}
